\chapter{Tổng quan và phạm vi}

\section{Mục tiêu}

RBAC Guard là nguyên mẫu kiểm soát truy cập theo vai trò kết hợp với phân tích
nhật ký bảo mật dựa trên luật. Dự án có hai bề mặt sử dụng bổ sung cho nhau.
Phần thứ nhất là pipeline Python nhận log CSV hoặc JSON, đối chiếu quyền trong
SQLite, tạo cảnh báo và xuất artifact có thể tái lập. Phần thứ hai là demo Nova
Bank gồm giao diện Next.js và API FastAPI, dùng để thể hiện xác thực, kiểm tra
quyền ở backend, phân tách nhiệm vụ và audit trail trong một luồng giao dịch
có thể trình diễn.

Mục tiêu kỹ thuật của hệ thống là trả lời ba câu hỏi vận hành: ai được phép làm
gì trên tài nguyên nào; sự kiện nào có dấu hiệu rủi ro; và bằng chứng nào giải
thích một quyết định từ chối hoặc một cảnh báo. Hệ thống ưu tiên tính minh bạch:
mỗi alert giữ rule, event nguồn, evidence, điểm rủi ro và severity thay vì chỉ
trả về một nhãn tổng quát.

\section{Phạm vi}

Pipeline phát hiện ba nhóm rủi ro nền: dữ liệu đầu vào mang dấu hiệu SQL
Injection, nhiều lần đăng nhập thất bại trong cửa sổ thời gian và truy cập
không được mô hình RBAC cấp quyền. Khi bật chế độ context-risk, hệ thống bổ
sung các tín hiệu hoạt động ngoài giờ, IP mới, tài nguyên hiếm và từ chối lặp
lại; các alert liên quan còn có thể được gom thành incident theo cặp người
dùng--IP và cửa sổ thời gian.

Demo Nova Bank tập trung vào bài toán phân tách nhiệm vụ trong phê duyệt giao
dịch. Ở chế độ kiểm soát cơ bản, một Teller có thể tạo và hoàn tất giao dịch.
Khi áp dụng chính sách RBAC, Teller bị backend từ chối khi truy cập nghiệp vụ
phê duyệt; chỉ Controller có quyền tương ứng. Các request được cho phép, từ
chối và thay đổi chính sách đều được ghi vào audit log. Adapter audit chuyển
các bản ghi xác thực và phân quyền này thành Event chuẩn hóa, nhờ đó quyết định
của Nova Bank có thể đi qua cùng pipeline phát hiện và báo cáo. Log request của
gateway vẫn là nguồn riêng cho SQL Injection vì audit nghiệp vụ không lưu toàn
bộ payload HTTP.

\section{Đóng góp kỹ thuật}

Báo cáo có bốn đóng góp có thể kiểm chứng. Thứ nhất, một mô hình RBAC có
enforcement ở API và policy separation of duties cho luồng phê duyệt giao dịch.
Thứ hai, một pipeline rule-based có hợp đồng Event, evidence, risk score và
artifact tái lập. Thứ ba, adapter nối audit log Nova Bank với pipeline cùng
thực nghiệm đầu-cuối đi từ thao tác nghiệp vụ tới alert. Thứ tư, phương pháp đánh
giá mở rộng gồm bộ regression, holdout 300 Event có manifest, metric theo Event
và campaign, cùng phân tích độ nhạy của luật password guessing.

\section{Ngoài phạm vi}

Đây là Proof of Concept dùng dữ liệu tổng hợp, SQLite cục bộ và luật xác định
trước. Nó không phải hệ thống ngân hàng triển khai thực tế, SIEM thời gian
thực hay cơ chế phát hiện dùng machine learning. Kết quả metric chứng minh
implementation khớp với scenario có nhãn của dự án. Bộ holdout làm giảm việc
đánh giá chỉ trên 60 Event ban đầu nhưng vẫn là dữ liệu tổng hợp do cùng nhóm
thiết kế, không đại diện cho hiệu năng trên nhật ký vận hành độc lập. Các giới
hạn này được giữ nguyên trong mọi diễn giải kết quả ở
Chương~\ref{chap:verification}.

\section{Tiêu chí hoàn thành}

Một lần chạy hợp lệ phải có thể khởi tạo dữ liệu RBAC, phân tích log, sinh
alert và đánh giá bằng đúng lệnh CLI được công bố. Với demo web, quyền phải
được kiểm tra bởi FastAPI thay vì chỉ ẩn nút ở giao diện; việc truy cập trực
tiếp URL phê duyệt của Teller sau khi áp dụng separation of duties phải trả về
HTTP 403 và không thay đổi trạng thái giao dịch. Toàn bộ các điều kiện này được
đối chiếu bằng kiểm tra tự động, thực nghiệm dữ liệu hoặc kịch bản trình diễn
tương ứng.
